% GENERATED FILE. Edit canonical lessons, then run npm run generate:books.
% canonical-source-hash: fnv1a64:f031539a2d1d675d
% canonical-lessons: ZH-C503-shufa, ZH-C503-yishu, ZH-C503-zuopin, ZH-C503-chuangzuo, ZH-C503-nianhua, ZH-R503-review-498-503

\chapter{Art}
\label{ch:zh-art}
\hlchaptermodality{503}

\begin{chapteropening}
\textbf{By the end of this chapter:} \emph{I can talk about calligraphy, the arts, a work of art, creating one and new year prints.}

Complete the last lesson of chapter 503: I can talk about calligraphy, the arts, a work of art, creating one and new year prints.
\end{chapteropening}

\section[\texorpdfstring{\zh{书法} (shūfǎ)}{书法 (shūfǎ)}]{\zh{书法} (shūfǎ) — calligraphy}

\label{lesson:ZH-C503-shufa}



\begin{quote}
Before the new one: say the Mandarin for a famous painting, then the Mandarin for an art exhibition.
\end{quote}

\subsection*{You'll want to know: \zh{书法}}
\textbf{\zh{书法}} — \emph{shūfǎ} — \textquotedblleft{}calligraphy\textquotedblright{}.

\begin{quote}
\emph{wǒ yéye xué le shí nián shūfǎ} — My grandpa studied calligraphy for ten years.
\end{quote}

\begin{cousinweb}
\textbf{\zh{书法}} (\emph{shūfǎ}) is built from \zh{书} 'writing' and \zh{法} 'a rule', both inherited from Middle Chinese, \emph{syo} and \emph{pjop} (Guangyun, in Baxter's transcription), which glosses the second 'a rule' and 'a norm'.
\end{cousinweb}

\begin{grammarlens}[title={What you've built}]
One more word for saying how it was.
\end{grammarlens}

\subsection*{Your turn}
\begin{itemize}
\raggedright
  \item \emph{Say it:} \emph{shūfǎ}
  \item \emph{Say it:} \emph{shūfǎ}, once more
  \item \emph{Say it:} \emph{wǒ yéye xué le shí nián shūfǎ} — My grandpa studied calligraphy for ten years.
  \item \emph{Recall:} say the Mandarin for a famous painting, then the Mandarin for an art exhibition, then say \emph{shūfǎ} again
\end{itemize}

\begin{tcolorbox}[breakable,colback=teal!4,colframe=teal!35!black,title={Before you move on}]
What does \zh{书法} mean? (\textquotedblleft{}Calligraphy\textquotedblright{}.) Say it once more.
\end{tcolorbox}
\section[yìshù]{yìshù — art; the arts}

\label{lesson:ZH-C503-yishu}



\begin{quote}
Before the new one: say the Mandarin for an art exhibition, then the Mandarin for calligraphy.
\end{quote}

\subsection*{You'll want to know: yìshù}
\emph{yìshù} — \textquotedblleft{}art; the arts\textquotedblright{}.

\begin{quote}
\emph{tā zuì xǐhuan yìshù} — Art is what she loves most.
\end{quote}

Its characters come later; for now it is a word for the ear.

\begin{cousinweb}
\emph{yìshù} is built from \emph{yì} 'a skill' and \emph{shù} 'a technique', both inherited from Middle Chinese, \emph{ngjiejH} and \emph{zywit} (Guangyun, in Baxter's transcription), which glosses the first 'ability' and the second 'technique'.
\end{cousinweb}

\begin{grammarlens}[title={What you've built}]
One more word for saying how it was.
\end{grammarlens}

\subsection*{Your turn}
\begin{itemize}
\raggedright
  \item \emph{Say it:} \emph{yìshù}
  \item \emph{Say it:} \emph{yìshù}, once more
  \item \emph{Say it:} \emph{tā zuì xǐhuan yìshù} — Art is what she loves most.
  \item \emph{Recall:} say the Mandarin for an art exhibition, then the Mandarin for calligraphy, then say \emph{yìshù} again
\end{itemize}

\begin{tcolorbox}[breakable,colback=teal!4,colframe=teal!35!black,title={Before you move on}]
What does yìshù mean? (\textquotedblleft{}Art; the arts\textquotedblright{}.) Say it once more.
\end{tcolorbox}
\section[zuòpǐn]{zuòpǐn — a work of art; a piece of writing}

\label{lesson:ZH-C503-zuopin}



\begin{quote}
Before the new one: say the Mandarin for calligraphy, then the Mandarin for the arts.
\end{quote}

\subsection*{You'll want to know: zuòpǐn}
\emph{zuòpǐn} — \textquotedblleft{}a work of art; a piece of writing\textquotedblright{}.

\begin{quote}
\emph{zhè shì tā zuì hǎo de zuòpǐn} — This is his best work.
\end{quote}

Its characters come later; for now it is a word for the ear.

\begin{cousinweb}
\emph{zuòpǐn} is built from \zh{作} 'to make' and \emph{pǐn} 'a kind', both inherited from Middle Chinese, \emph{tsak} and \emph{phimX} (Guangyun, in Baxter's transcription), which glosses the second 'a rank' and 'a kind'.
\end{cousinweb}

\begin{grammarlens}[title={What you've built}]
One more word for saying how it was.
\end{grammarlens}

\subsection*{Your turn}
\begin{itemize}
\raggedright
  \item \emph{Say it:} \emph{zuòpǐn}
  \item \emph{Say it:} \emph{zuòpǐn}, once more
  \item \emph{Say it:} \emph{zhè shì tā zuì hǎo de zuòpǐn} — This is his best work.
  \item \emph{Recall:} say the Mandarin for calligraphy, then the Mandarin for the arts, then say \emph{zuòpǐn} again
\end{itemize}

\begin{tcolorbox}[breakable,colback=teal!4,colframe=teal!35!black,title={Before you move on}]
What does zuòpǐn mean? (\textquotedblleft{}A work of art; a piece of writing\textquotedblright{}.) Say it once more.
\end{tcolorbox}
\section[chuàngzuò]{chuàngzuò — to create (art or writing)}

\label{lesson:ZH-C503-chuangzuo}



\begin{quote}
Before the new one: say the Mandarin for the arts, then the Mandarin for a work of art.
\end{quote}

\subsection*{You'll want to know: chuàngzuò}
\emph{chuàngzuò} — \textquotedblleft{}to create (art or writing)\textquotedblright{}.

\begin{quote}
\emph{tā yòng le wǔ nián chuàngzuò zhège zuòpǐn} — He spent five years creating this work.
\end{quote}

Its characters come later; for now it is a word for the ear.

\begin{cousinweb}
\emph{chuàngzuò} is built from \emph{chuàng} 'to found' and \zh{作} 'to make', both inherited from Middle Chinese, \emph{tsrhjangH} and \emph{tsak} (Guangyun, in Baxter's transcription). The dictionary glosses the level-tone reading 'a wound', and gives the falling-tone one as another spelling of a word glossed 'to begin' and 'to found a craft'.
\end{cousinweb}

\begin{grammarlens}[title={What you've built}]
One more word for saying how it was.
\end{grammarlens}

\subsection*{Your turn}
\begin{itemize}
\raggedright
  \item \emph{Say it:} \emph{chuàngzuò}
  \item \emph{Say it:} \emph{chuàngzuò}, once more
  \item \emph{Say it:} \emph{tā yòng le wǔ nián chuàngzuò zhège zuòpǐn} — He spent five years creating this work.
  \item \emph{Recall:} say the Mandarin for the arts, then the Mandarin for a work of art, then say \emph{chuàngzuò} again
\end{itemize}

\begin{tcolorbox}[breakable,colback=teal!4,colframe=teal!35!black,title={Before you move on}]
What does chuàngzuò mean? (\textquotedblleft{}To create (art or writing)\textquotedblright{}.) Say it once more.
\end{tcolorbox}
\section[\texorpdfstring{\zh{年画} (niánhuà)}{年画 (niánhuà)}]{\zh{年画} (niánhuà) — a new year print (hung at the spring festival)}

\label{lesson:ZH-C503-nianhua}



\begin{quote}
Before the new one: say the Mandarin for a work of art, then the Mandarin for to create (art or writing).
\end{quote}

\subsection*{You'll want to know: \zh{年画}}
\textbf{\zh{年画}} — \emph{niánhuà} — \textquotedblleft{}a new year print (hung at the spring festival)\textquotedblright{}.

\begin{quote}
\emph{guònián de shíhou wǒ nǎinai mǎi le yìxiē niánhuà} — At the new year, my grandma bought some new year prints.
\end{quote}

\begin{cousinweb}
\textbf{\zh{年画}} (\emph{niánhuà}) is built from \zh{年} 'a year' and \zh{画} 'a painting', both inherited from Middle Chinese, \emph{nen} and \emph{hweaH} (Guangyun, in Baxter's transcription), which glosses the first 'the grain is ripe'.
\end{cousinweb}

\begin{grammarlens}[title={What you've built}]
One more word for saying how it was.
\end{grammarlens}

\subsection*{Your turn}
\begin{itemize}
\raggedright
  \item \emph{Say it:} \emph{niánhuà}
  \item \emph{Say it:} \emph{niánhuà}, once more
  \item \emph{Say it:} \emph{guònián de shíhou wǒ nǎinai mǎi le yìxiē niánhuà} — At the new year, my grandma bought some new year prints.
  \item \emph{Recall:} say the Mandarin for a work of art, then the Mandarin for to create (art or writing), then say \emph{niánhuà} again
\end{itemize}

\begin{tcolorbox}[breakable,colback=teal!4,colframe=teal!35!black,title={Before you move on}]
What does \zh{年画} mean? (\textquotedblleft{}A new year print (hung at the spring festival)\textquotedblright{}.) Say it once more.
\end{tcolorbox}
\section[(dialogue)]{Looking back: from On the Pitch to Art}

\label{lesson:ZH-R503-review-498-503}



\begin{quote}
Say the Mandarin for physical education and a sports pitch.
\end{quote}

\subsection*{Your turn}
\begin{itemize}
\raggedright
  \item \emph{Say it:} physical education, a sports pitch, the goalmouth, to shoot at goal, to score a goal
  \item \emph{Say it:} a sports team, a player in a ball game, a team captain, to do training (for a sport), to pass the ball
  \item \emph{Say it:} to win a contest, a ball game, final (of a competition), the score of a game, to take the field
  \item \emph{Say it:} an athlete, the champion, a gold medal, to win (a prize), a star player
  \item \emph{Say it:} an oil painting, a traditional Chinese painting, a painter (an artist), a famous painting, an art exhibition
  \item \emph{Say it:} calligraphy, the arts, a work of art, to create (art or writing), a new year print
  \item \emph{Read it:} \zh{球} \zh{赛} \zh{画} — each one aloud, then a word that holds it
\end{itemize}

\begin{tcolorbox}[breakable,colback=teal!4,colframe=teal!35!black,title={Before you move on}]
Physical education? (\emph{tǐyù}.) A sports pitch? (\textbf{\zh{球场}}, \emph{qiúchǎng}.) The goalmouth? (\textbf{\zh{球门}}, \emph{qiúmén}.) To shoot at goal? (\emph{shèmén}.) To score a goal? (\textbf{\zh{进球}}, \emph{jìnqiú}.) A sports team? (\emph{qiúduì}.) A player in a ball game? (\emph{qiúyuán}.) A team captain? (\emph{duìzhǎng}.) To do training (for a sport)? (\emph{xùnliàn}.) To pass the ball? (\textbf{\zh{传球}}, \emph{chuánqiú}.) To win a contest? (\emph{huòshèng}.) A ball game? (\textbf{\zh{球赛}}, \emph{qiúsài}.) Final (of a competition)? (\emph{juésài}.) The score of a game? (\emph{bǐfēn}.) To take the field? (\textbf{\zh{上场}}, \emph{shàngchǎng}.) An athlete? (\emph{yùndòngyuán}.) The champion? (\emph{guànjūn}.) A gold medal? (\emph{jīnpái}.) To win (a prize)? (\emph{huòdé}.) A star player? (\textbf{\zh{球星}}, \emph{qiúxīng}.) An oil painting? (\emph{yóuhuà}.) A traditional Chinese painting? (\textbf{\zh{国画}}, \emph{guóhuà}.) A painter (an artist)? (\textbf{\zh{画家}}, \emph{huàjiā}.) A famous painting? (\textbf{\zh{名画}}, \emph{mínghuà}.) An art exhibition? (\emph{huàzhǎn}.) Calligraphy? (\textbf{\zh{书法}}, \emph{shūfǎ}.) The arts? (\emph{yìshù}.) A work of art? (\emph{zuòpǐn}.) To create (art or writing)? (\emph{chuàngzuò}.) A new year print? (\textbf{\zh{年画}}, \emph{niánhuà}.)
\end{tcolorbox}
